\subsection{Явный вид ограничений и условие оптимальности}
\label{subsec:transport-duality-anchor}

Запишем в явном виде ограничения закрытой задачи:
\[
  \underbrace{
  \begin{cases}
    x_{11} + x_{12} + \dots + x_{1n} = a_1, \\
    x_{21} + x_{22} + \dots + x_{2n} = a_2, \\
    \quad \dots \\
    x_{m1} + x_{m2} + \dots + x_{mn} = a_m
  \end{cases}}_{m \ \text{«строчных» равенств}}
  \qquad
  \underbrace{
  \begin{cases}
    x_{11} + x_{21} + \dots + x_{m1} = b_1, \\
    x_{12} + x_{22} + \dots + x_{m2} = b_2, \\
    \quad \dots \\
    x_{1n} + x_{2n} + \dots + x_{mn} = b_n
  \end{cases}}_{n \ \text{«столбцовых» равенств}}
\]

Если переменные $x_{ij}$ вытянуть в один вектор $x[N]$ по строкам
($N$ — $mn$ индексов), то матрица ограничений имеет блочный вид
\[
  A[M,N] =
  \begin{pmatrix}
    \mathbf{1}\T & 0 & \cdots & 0 \\
    0 & \mathbf{1}\T & \cdots & 0 \\
    \vdots & \vdots & \ddots & \vdots \\
    0 & 0 & \cdots & \mathbf{1}\T \\[4pt]
    E & E & \cdots & E
  \end{pmatrix},
  \qquad
  \begin{aligned}
    &\mathbf{1}\T = (1\ 1\ \dots\ 1) \in \R^{n}, \\
    &E \ \text{— единичная } n \times n ,
  \end{aligned}
\]
размера $(m+n) \times mn$: первые $m$ строк — «строчные» равенства (по $n$ единиц
подряд), последние $n$ строк — «столбцовые» (единичные блоки). При этом
\[
  \sum_{i=1}^{m} a_i = \sum_{j=1}^{n} b_j = P .
\]

\medskip
\noindent\textbf{Условие оптимальности:}
\[
  c\T[N] - y\T[M]\, A[M,N] \ge 0 ,
  \qquad
  \big(c\T[N] - y\T[M]\, A[M,N]\big)\, x[N] = 0 .
\]

\begin{idea}[title={Откуда берутся $v_j - u_i \le c_{ij}$ и $v_j - u_i = c_{ij}$}]
Это и есть ответ на вопросительные знаки в тетради. Условие оптимальности выше —
стандартная пара «двойственная допустимость $+$ дополняющая нежёсткость»; надо
лишь расписать её покомпонентно.

\medskip
\textbf{Шаг 1: что такое $y[M]$.} Вектор $y[M]$ — двойственные переменные, по одной
на каждое ограничение. Ограничений $m + n$: первые $m$ отвечают строкам (пунктам
хранения), последние $n$ — столбцам (пунктам назначения). Обозначим их
\[
  y[M] = (\,\underbrace{y_1, \dots, y_m}_{\text{строки}},\
            \underbrace{y_{m+1}, \dots, y_{m+n}}_{\text{столбцы}}\,) .
\]

\medskip
\textbf{Шаг 2: как выглядит столбец $A$.} Переменная $x_{ij}$ входит ровно в два
ограничения: в $i$-е «строчное» и в $j$-е «столбцовое». Значит её столбец в
матрице $A$ имеет две единицы — в позициях $i$ и $m+j$ — и нули в остальных.
Поэтому компонента произведения $y\T A$, отвечающая $x_{ij}$, равна
\[
  \big(y\T A\big)_{ij} = y_i + y_{m+j} .
\]

\medskip
\textbf{Шаг 3: покомпонентная запись.} Первое условие
$c\T - y\T A \ge 0$ означает
\[
  c_{ij} - \big(y_i + y_{m+j}\big) \ge 0 \qquad \forall i, j .
\]
Положим
\[
  u_i = -\,y_i , \qquad v_j = y_{m+j} .
\]
Тогда $y_i + y_{m+j} = v_j - u_i$, и условие превращается ровно в
\[
  \boxed{\,v_j - u_i \le c_{ij} \quad \forall i, j\,}
\]
— то самое неравенство из пункта~2 алгоритма. Так потенциалы $u_i$ и $v_j$
оказываются просто двойственными переменными транспортной задачи, взятыми с
удобными знаками.

\medskip
\textbf{Шаг 4: вторая строка — дополняющая нежёсткость.} Произведение
$\big(c\T - y\T A\big) x$ — это сумма
\[
  \sum_{i,j} \underbrace{\big(c_{ij} - (v_j - u_i)\big)}_{\ge\, 0}\,
             \underbrace{x_{ij}}_{\ge\, 0} = 0 .
\]
Сумма неотрицательных слагаемых равна нулю только если каждое слагаемое нулевое.
Значит для всякой клетки с $x_{ij} > 0$ (то есть \emph{заполненной}) обязано
выполняться
\[
  \boxed{\,v_j - u_i = c_{ij}\,}
\]
— а это и есть система, из которой в пункте~1 алгоритма находят потенциалы.

\medskip
\textbf{Итог.} Две строки условия оптимальности распадаются ровно на два шага
метода потенциалов: равенства по заполненным клеткам дают $u_i$ и $v_j$,
неравенства по пустым проверяют оптимальность. Если хоть в одной пустой клетке
$v_j - u_i > c_{ij}$, двойственная допустимость нарушена, план не оптимален, и
эту клетку вводят в базис циклом пересчёта.

Система на потенциалы недоопределена на одну переменную (уравнений $m+n-1$,
неизвестных $m+n$) именно потому, что строки матрицы $A$ линейно зависимы:
сумма первых $m$ строк равна сумме последних $n$. Поэтому один потенциал
назначают произвольно — например $v_1 = 0$.
\end{idea}

\subsection{Разрешимость транспортной задачи}

\begin{theorem}[о транспортной задаче]\label{th:transport-solvable}
Любая закрытая транспортная задача имеет решение.
\end{theorem}

\begin{proof}
Допустимое множество непусто: положим для всех $i, j$
\[
  x_{ij} := \frac{a_i\, b_j}{P} .
\]
Подстановка этих значений показывает допустимость вектора:
\[
  \sum_{j=1}^{n} x_{ij} = \frac{a_i}{P} \sum_{j=1}^{n} b_j = \frac{a_i P}{P} = a_i ,
  \qquad
  \sum_{i=1}^{m} x_{ij} = \frac{b_j}{P} \sum_{i=1}^{m} a_i = b_j ,
\]
и очевидно $x_{ij} \ge 0$.

Целевая функция ограничена сверху и снизу. Действительно, существуют
\[
  \underline{c} = \min_{i,j} c_{ij} , \qquad \overline{c} = \max_{i,j} c_{ij} ,
\]
а суммарный объём перевозок у любого допустимого плана один и тот же:
$\sum_{i,j} x_{ij} = P$. Отсюда
\[
  \underline{c}\,P \ \le\ \sum_{i}\sum_{j} c_{ij} x_{ij} \ \le\ \overline{c}\,P .
\]
Непустое допустимое множество плюс ограниченная снизу целевая функция дают
существование оптимального решения.
\end{proof}

\begin{idea}
Строго говоря, ограниченности снизу мало — нужна ещё замкнутость допустимого
множества и достижимость минимума. Здесь это выполнено: допустимое множество
закрытой задачи — многогранное и \emph{ограниченное} ($0 \le x_{ij} \le \min\{a_i, b_j\}$),
то есть компакт, а линейная функция на компакте достигает минимума. Иначе это же
следует из общей теории: по теореме~\ref{th:optimal-is-support} оптимум достигается
в одном из конечного числа опорных векторов.
\end{idea}

Для приближения классической постановки транспортной задачи к реальным
рассматриваются усложнения транспортной задачи.
